% =========================================================================
% AI 工具使用声明（置于参考文献之前）
% -------------------------------------------------------------------------
% 按 2026 年竞赛 AI 工具使用规定，如实说明工具、使用环节、提示过程与人工核验。
% 【】内的内容需由队员按实际情况填写，不要在提交前保留占位符。
% =========================================================================

\section*{人工智能工具使用声明}

本队在本次竞赛的建模与写作过程中使用了人工智能辅助工具。为便于评阅核查，现如实说明使用情况如下。

\textbf{（一）所用工具。}【待填：逐项写明工具名称与具体型号或版本号，例如 ``XXX（版本号/日期）''；若使用了多个工具，分别列出并说明各自用于哪些环节。】

\textbf{（二）使用环节。}本次使用覆盖以下环节：
\begin{enumerate}
  \item 模型思路的讨论与比较（例如在确定性调度、序贯决策、滚动调整与跨日规划之间取舍）；
  \item 公式与符号的整理与一致性核对（例如分位裕量的定义、双向违约价的分段结算方式）；
  \item Python 程序的编写、调试与重构（线性规划装配、执行层仿真、结果工作簿写出）；
  \item 代码审查与数值审计（读代码核对论文公式，独立复算费用账本与边界残差）；
  \item 论文文字与排版的校对（术语统一、交叉引用、图表编号与版式检查）。
\end{enumerate}

\textbf{（三）主要提示过程。}【待填：摘录若干条能反映工作方式的代表性提示，例如``请核对论文中分位数裕量公式与代码实现对同一窗口残差的定义是否一致''``请检查结算式中最终合同费与调整附加费是否存在重复计费''。不必逐条罗列全部对话。】

\textbf{（四）人工采纳与核验。}【待填：说明哪些建议被采纳、哪些被否决及其理由；并说明关键数值的核验方式，例如``论文引用的全部费用数字由脚本重新运行输出，并与独立复算路径逐位比对''``三张指定日期结果表由程序直接生成，未手工录入''。】

\textbf{（五）责任声明。}本文的模型设定、参数取值、全部数值结果与最终文字由队员确定并负责。AI 工具产出的内容均经队员逐项核对；凡涉及数据、公式与计算的部分，均以本文脚本的实际运行输出为准。题目要求的原创性工作（模型结构设计、对照实验设计、结果解释）由队员完成。
